\documentclass[aspectratio=169]{beamer}

\usetheme{Madrid}

\usepackage{amsmath}
\usepackage{listings}
\usepackage{xcolor}

% Code display settings
\lstset{
basicstyle=\ttfamily\small,
frame=single,
breaklines=true
}

\title{Time Series}
\subtitle{Decomposition: Seasonality Estimation}
\author{Mr. Aymen Ben Brik}
\institute{Time Series Course 1GAMMA}
\date{\today}

\begin{document}

%------------------------------------------------
\begin{frame}
\titlepage
\end{frame}

%------------------------------------------------
\begin{frame}{Course Objectives}

By the end of this session, students will be able to:

\begin{itemize}

\item Understand the concept of seasonality in time series

\item Estimate seasonal components using different methods

\item Apply seasonal averaging

\item Apply parametric regression for seasonality estimation

\item Apply cosine-sine models

\item Remove seasonality from time series data

\end{itemize}

\end{frame}

%------------------------------------------------
\begin{frame}{Time Series Model}

Time series decomposition:

\[
Y_t = T_t + S_t + \varepsilon_t
\]

Where:

\begin{itemize}

\item $T_t$ : trend

\item $S_t$ : seasonality

\item $\varepsilon_t$ : stationary component

\end{itemize}

Goal:

\begin{itemize}
\item Estimate seasonality
\item Remove seasonality
\end{itemize}

\end{frame}

%------------------------------------------------
\begin{frame}{Seasonal Average Method}

Seasonal component estimation:

\[
\hat{S}_k =
\frac{1}{n_k}
\sum Y_t
\]

where observations belong to season $k$

Deseasonalized series:

\[
Y_t - \hat{S}_k
\]

\end{frame}

%------------------------------------------------
\begin{frame}[fragile]{Example: Seasonal Average in R}

\begin{lstlisting}[language=R]

# Monthly data example
month <- rep(1:12,10)
y <- 10 + sin(month/12*2*pi) + rnorm(120)

# Estimate seasonal means
season <- tapply(y, month, mean)

# Remove seasonality
y_deseason <- y - season[month]

\end{lstlisting}

\end{frame}

%------------------------------------------------
\begin{frame}[fragile]{Example: Seasonal Average in Python}

\begin{lstlisting}

import numpy as np
import pandas as pd

month = np.tile(np.arange(1,13),10)
y = 10 + np.sin(month/12*2*np.pi) + np.random.randn(120)

df = pd.DataFrame({"y":y,"month":month})

season = df.groupby("month").mean()

df["deseason"] = df["y"] - df["month"].map(season["y"])

\end{lstlisting}

\end{frame}

%------------------------------------------------
\begin{frame}{Parametric Regression Method}

Model using dummy variables:

\[
Y_t =
\beta_0 +
\beta_1 D_1 +
\dots +
\beta_d D_d +
\varepsilon_t
\]

Where dummy variables represent seasons.

\end{frame}

%------------------------------------------------
\begin{frame}[fragile]{Example: Seasonal Regression in R}

\begin{lstlisting}[language=R]

month <- factor(month)

model <- lm(y ~ month)

season <- predict(model)

\end{lstlisting}

\end{frame}

%------------------------------------------------
\begin{frame}[fragile]{Example: Seasonal Regression in Python}

\begin{lstlisting}

import statsmodels.formula.api as smf

df["month"] = df["month"].astype("category")

model = smf.ols("y ~ C(month)", data=df).fit()

season = model.predict(df)

\end{lstlisting}

\end{frame}

%------------------------------------------------
\begin{frame}{Cosine-Sine Model}

Seasonality modeled as:

\[
S_t =
a \cos(\omega t)
+
b \sin(\omega t)
\]

Where:

\[
\omega =
\frac{2\pi}{d}
\]

\end{frame}

%------------------------------------------------
\begin{frame}[fragile]{Example: Cosine-Sine Model in R}

\begin{lstlisting}[language=R]

omega <- 2*pi/12

model <- lm(y ~ cos(omega*t) + sin(omega*t))

season <- predict(model)

\end{lstlisting}

\end{frame}

%------------------------------------------------
\begin{frame}[fragile]{Example: Cosine-Sine Model in Python}

\begin{lstlisting}

omega = 2*np.pi/12

df["cos"] = np.cos(omega*t)
df["sin"] = np.sin(omega*t)

import statsmodels.api as sm

X = df[["cos","sin"]]
X = sm.add_constant(X)

model = sm.OLS(y,X).fit()

season = model.predict(X)

\end{lstlisting}

\end{frame}

%------------------------------------------------
\begin{frame}{Deseasonalization}

Remove seasonality:

\[
Y_t^{*} =
Y_t -
\hat{S}_t
\]

Result:

\begin{itemize}

\item Stationary series

\item Ready for time series modeling

\end{itemize}

\end{frame}

%------------------------------------------------
\begin{frame}{Summary}

Methods covered:

\begin{itemize}

\item Seasonal averaging

\item Regression with dummy variables

\item Cosine-sine model

\end{itemize}

\end{frame}

%------------------------------------------------
\end{document}